\subsection{The auxiliary-circle limit}
\label{subsec:x9-circle-limit}

The marked four-node locus has a large-radius limit on the original
five-dimensional contraction face. Its four-hypermultiplet infrared
description, however, has no radius-independent range of validity.
These are distinct statements: the limiting geometry can be identified
without treating the finite-circle conifold theory as the full
five-dimensional theory.

\paragraph{Suppression of smooth-side instantons.}
The smooth-side effective cone gives a further all-orders consequence. By
\eqref{eq:x9-neutral-common-mori}, every nonzero effective degree
has $m+n>0$. Hence $J_0=\alpha K_3+\beta K_2$, with
$\alpha,\beta>0$, lies in the interior of the smooth-side K\"ahler
cone: its degree on a curve is $\alpha n+\beta m$.
The ambient boundary curves of degree $(1,0,0)$ do not meet the
Calabi--Yau hypersurface. Along
\eqref{eq:x9-circle-scaling-path}, every smooth-side instanton is thus
suppressed, although $v=v_0$ remains fixed. The convergent
Frobenius series and the locally invertible mirror map give an
analytic nonpolynomial prepotential near the origin. To see the
convergence directly, the factorial ratios and their first two
derivatives are bounded by an exponential in total degree times
a polynomial in that degree; the latter includes the harmonic-number
factors. Analytic inversion of the mirror map and integration of the
regular magnetic series preserve a nonzero convergence domain. Its support
has no pure $q_1$ term, so, for any fixed number of flat-coordinate
derivatives,
\begin{equation}
 \partial_{\tau_{i_1}}\cdots\partial_{\tau_{i_k}}
 F^{\rm sm}_{\rm inst}
 =O\!\left(e^{-2\pi L\min(\alpha,\beta)}\right).
 \label{eq:x9-neutral-instanton-suppression}
\end{equation}
The bound holds at fixed sufficiently small $v_0$ on the smooth-side
large-volume branch. Consequently its quantum vector couplings have
the cubic five-dimensional limit already matched in
\eqref{eq:x9-primitive-cubic}.
\paragraph{Radius and period normalization.}
Let $R_{\rm aux}$ be the radius of the spacetime circle reducing five
dimensions to four, measured in the five-dimensional metric, and put
$\mathcal R=R_{\rm aux}M$ for a fixed mass unit $M$.
This is not the elliptic-fiber circle of Section~\ref{sec:fibrations}.
In the leading circle dictionary the dimensionless vector coordinates
are $t_a=B_a+i\mathcal R h_a$, where the real five-dimensional scalars
obey $\frac16\kappa_{abc}h_ah_bh_c=1$ after volume normalization
\cite[eqs.~(5.35)--(5.41)]{Hattab:2025M2}.
The total M-theory Calabi--Yau volume is a separate hypermultiplet
parameter and is kept fixed and finite here. Calibrating the periods
by the unit D0 charge gives, for an existing BPS state,
\begin{equation}
 M_\gamma^{(5\text{-metric})}
       =\frac{|Z_\gamma|}{R_{\rm aux}}.
 \label{eq:x9-radius-mass-normalization}
\end{equation}
The four-dimensional Einstein-frame mass has an additional common Weyl
factor; it is not the mass used in \eqref{eq:x9-radius-mass-normalization}.
The same KK normalization is used for local circle theories in
\cite[eq.~(2.27)]{Closset:2022UPlane}.

\paragraph{A scaling inside the marked family.}
Fix sufficiently small positive $v_0,a_0,b_0$ and constants
$\alpha,\beta>0$. With $L\to\infty$, take
\begin{equation}
 z_1=z_5=z_6=1,\qquad v=v_0,\qquad
 a=a_0e^{-2\pi L\beta},\qquad b=b_0e^{-2\pi L\alpha}.
 \label{eq:x9-circle-scaling-path}
\end{equation}
Choose generic constants so that the path lies in the four-node
stratum for all sufficiently large $L$. This path stays in the
convergence domain of the marked period resummation. Writing
$t_j=iy_j$ for $j=2,3,4$, its exact flat coordinates satisfy
\begin{equation}
 \begin{aligned}
 t_1=t_5=t_6&=0,&
 y_2&=c+O(a+b),\qquad c=-\frac{\log v_0}{2\pi}>0,\\
 y_3&=L\beta+O(1),&y_4&=L\alpha+O(1).
 \end{aligned}
 \label{eq:x9-circle-period-scaling}
\end{equation}
Indeed all logarithmic corrections are bounded analytic functions
of $(v,a,b)$ by \eqref{eq:x9-bulk-electric-periods}.
At $a=b=0$, \eqref{eq:x9-local-exact-endpoint} gives $q_2=v$,
which fixes the first limit in \eqref{eq:x9-circle-period-scaling}.
Thus, for $J_\Sigma=y_2H_2+y_3H_3+y_4H_4$,
\begin{equation}
 \frac{J_\Sigma}{L}\longrightarrow
 J_0=\alpha D_0+\beta D_1,\qquad
 \frac{J_\Sigma^3}{6L^3}\longrightarrow
 \mathcal V_0=
 \frac{83\alpha^3+147\alpha^2\beta+81\alpha\beta^2+14\beta^3}{6}>0.
 \label{eq:x9-circle-geometric-limit}
\end{equation}
The leading radius therefore scales as
$\mathcal R\sim L\mathcal V_0^{1/3}$, and the volume-normalized
five-dimensional class tends to $J_0/\mathcal V_0^{1/3}$.
Bounded resummed magnetic corrections do not alter this leading
cubic scaling. We do not assume that the local instanton corrections
themselves vanish at the conifold locus.

The extra condition needed to contract the whole del Pezzo surface
is visible in its restriction lattice. In the basis $(h,e_1,e_2)$
of $E=\operatorname{Bl}_2\mathbb P^2$,
\begin{equation}
 H_2|_E=h,\qquad H_3|_E=H_4|_E=0,\qquad
 (D_8\cap E)=\beta_2.
 \label{eq:x9-circle-surface-restriction}
\end{equation}
The two exceptional curves have zero period on $\Sigma$, whereas
$\beta_2$ has period $iy_2\ne0$. Its rescaled area $y_2/L$ tends
to zero, so $J_\Sigma|_E/L\to0$: the entire surface, not only the
two exceptional curves, contracts. This is precisely the additional
condition distinguishing the five-dimensional face
\eqref{eq:local-scale-plane} from the codimension-three complex
four-node locus.

\paragraph{Finite bulk gaugings and the local strong-coupling limit.}
The limit in \eqref{eq:x9-circle-geometric-limit} does not send the
M-theory bulk volume to infinity. In particular it is not the rigid
limit of Section~\ref{sec:decoupling}. With fixed five-dimensional
Planck mass, dimensional reduction instead gives
$M_4^2=2\pi R_{\rm aux}M_5^3$ in the common five-dimensional metric.
The two vectors imposing cooperation remain normalizable.
This can be checked after removing the graviphoton from the Hodge metric:
\begin{equation}
 \mathcal G_J^{\rm vec}(P,Q)
 =-JPQ+\frac{(J^2P)(J^2Q)}{6\mathcal V_J}.
 \label{eq:x9-circle-projected-metric}
\end{equation}
Up to the common finite volume normalization, the block on
$(V_1,V_2)=(-D_6,D_5)$ at $J_0$ is
\begin{equation}
 \mathcal G_{J_0}^{\rm vec}|_{V_1,V_2}
 =\begin{pmatrix}d&\tau-d\\\tau-d&d\end{pmatrix},\qquad
 d=\frac{2(7\alpha^2+8\alpha\beta+2\beta^2)^2}{3\mathcal V_0},
 \quad\tau=3\alpha+2\beta.
 \label{eq:x9-circle-finite-gaugings}
\end{equation}
Its eigenvalues are $\tau$ and
\begin{equation}
 2d-\tau=
 \frac{143\alpha^4+289\alpha^3\beta+199\alpha^2\beta^2
                       +52\alpha\beta^3+4\beta^4}
      {83\alpha^3+147\alpha^2\beta+81\alpha\beta^2+14\beta^3},
 \label{eq:x9-circle-positive-eigenvalue}
\end{equation}
both finite and strictly positive. By contrast, put
$\eta=y_2/L$ and $J=J_0+\eta H_2$. The restriction data give
\begin{equation}
 \frac12J^2E=\frac{\eta^2}{2},\qquad
 \mathcal G_J^{\rm vec}(E,E)
       =3\eta+\frac{\eta^4}{6\mathcal V_J}\longrightarrow0.
 \label{eq:x9-circle-local-strong-coupling}
\end{equation}
Bounded $H_3,H_4$ offsets in \eqref{eq:x9-circle-period-scaling}
do not change these limits. Thus the del Pezzo string tension scale
vanishes and its Coulomb direction becomes strongly coupled, while
the two bulk gaugings remain dynamical. These are limiting
five-dimensional geometric couplings, not the four-dimensional
threshold couplings obtained by integrating out massless conifold
hypermultiplets.

\paragraph{Why four hypermultiplets do not suffice in the limit.}
The KK gap is $1/R_{\rm aux}$. At any fixed positive energy cutoff,
the number of Fourier modes of even one five-dimensional massless
field grows without bound. Hence the four-dimensional truncation
requires energies much smaller than $1/R_{\rm aux}$ and has no
uniform cutoff as $R_{\rm aux}\to\infty$.

Equations \eqref{eq:x9-circle-scaling-path}--\eqref{eq:x9-circle-geometric-limit}
identify the limiting Coulomb-branch locus with the original
five-dimensional cooperative contraction. They do not construct
the interacting five-dimensional Higgs action or its finite-Planck
metric. The choice of mirror crepant chamber and magnetic BPS
survival also remain separate questions.
